% source: https://github.com/pfradin/luadraw/discussions/337#discussioncomment-18501273 (pfradin, block 1)
\documentclass[border=5pt]{standalone}% compile with lualatex only
\usepackage[svgnames]{xcolor}
\usepackage[3d]{luadraw}%https://github.com/pfradin/luadraw
%https://ctan.org/pkg/luadraw
\usepackage{fourier-otf}
\begin{luacode*}
local ld = luadraw
local pt3d = ld.pt3d
function ld.cutfacet2(L,f,sg)
-- L is a list of facets
-- f = (x,y,z) -> f(x,y,z) in R
-- sg = ">" or "<"
-- the part satisfying f>0 or f<0 is kept
-- returns two lists of facets, the first one satisfying f>0 if sg=">" or f<0 if sg="<"
    if pt3d.isPoint3d(L[1]) then L = {L} end
    if L.vertices ~= nil then -- polyèdre
        L = ld.poly2facet(L)
    end
    local function test(A)
        local r = f(A.x,A.y,A.z)
        if sg== ">" then return r
        else return -r
        end
    end
    local function solve(A,fa,B,fb)
        local a, b, c, fc = A, B
        k = 0
        while (pt3d.abs2(b-a)>1e-12) or (k<=50) do
            c = (a+b)/2; fc = test(c)
            if fa*fc < 0 then b = c
            else a = c; fa = fc
            end
            k = k + 1 
        end
        return c
    end
    local res, res2 = {}, {}
    for _,F in ipairs(L) do
        local nb, aux, aux2 = #F, {}, {}
        local A1, B1 = nil, F[1]
        local p1, p2, I = nil, test(B1)
        if math.abs(p2)<1e-8 then p2 = 0 end
        for k = 2, nb+1 do
            if k == nb+1 then k = 1 end -- on ferme la facette
            A1 = B1; p1 = p2; B1 = F[k]; p2 = test(B1)
            if math.abs(p2) < 1e-8 then p2 = 0 end
            if (p1*p2 < 0) or (p2 == 0) then
                if p2 == 0 then I = B1 else I = solve(A1,p1,B1,p2) end
                if I ~= nil then 
                    table.insert(aux,I) ; table.insert(aux2,I)
                end
            end
            if (p2 > 0) and (p1 ~= nil) then table.insert(aux,B1) end
            if (p2 < 0) and (p1 ~= nil) then table.insert(aux2,B1) end
        end
        if #aux>2 then table.insert(res,aux) end
        if #aux2>2 then table.insert(res2,aux2) end
    end
    return res, res2 -- returns two list of facets
end
\end{luacode*}

\begin{document}
\begin{luadraw}{name=sphere_cut_by_tube}
local ld = luadraw
local pt3d = ld.pt3d
local Origin, vecI, vecJ, vecK, M = pt3d.Origin, pt3d.vecI, pt3d.vecJ, pt3d.vecK, pt3d.M
local cos, sin, pi = math.cos, math.sin, math.pi
local g = ld.graph3d:new{ window={-4,5,-6,4}, viewdir={-30,65}, size={12,12}}
require 'luadraw_spherical'
local r, R = 3, 4
local r2 = r^2
local f = function(x,y,z) return pt3d.abs2(M(x,y,z))-r2 end
local L = ld.parametric3d( function(t) return R*vecI+R*cos(t)*vecI+3*sin(t)*vecK end, -pi,pi,35,false,0)[1]
local T = ld.line2tube(L, 1, {nbfacet=20, close=true})
local T1, T2 = ld.cutfacet2(T,f,">")
local B = ld.border(T2)
local I, r1, n = table.unpack( g:Sphere_outline(Origin,r).data )
T1,T3 = ld.cutfacet(T1, {I,n})

g:Dfacet(T2, {color="gray", opacity=0.6, mode=ld.mShadedOnly}) --inside sphere
g:Dfacet(T3, {color="gray", opacity=0.6, mode=ld.mShadedOnly}) --outside sphere but hidden
g:Define_sphere({radius=r, opacity=0.7})
g:DSregion(B[1], {fill="DarkGreen", fillopacity=0.7, color="red", width=8}) -- intersection1
g:DSregion(B[2], {fill="DarkGreen", fillopacity=0.7, color="red", width=8}) -- intersection2
g:Dspherical()
g:Dfacet(T1, {color="gray", opacity=0.7, mode=ld.mShadedOnly}) -- outside
g:Show()
\end{luadraw}
\end{document}
